\section{現代テスト理論}
\subsection{授業内容}
\subsubsection{科目の中でのこのコマの位置づけ}
因子分析によって項目と被験者の特徴を分離して考えることができるようになった。ここで学力テストに目を向けると，単因子でよいことと従属変数がバイナリになっていることがわかる。
この特殊な測定法についてのモデルを考えるために，まずは通過率の概念を導入したうえで，累積正規分布とその近似としてのロジスティック曲線，および1,2,3PLモデルを紹介する。
これらのテストは新しいテスト理論とよばれるが，それはこれまでのテスト理論に含まれていた集団に依存した測定であったこと，完全データに限定されていたことなどを乗り越えられるからである。もちろんテストの等価がしやすいという側面もある。
テスト理論の展開としての項目反応理論と，因子分析モデルとの相同性を強調することで，見えないものを測定しようとするアプローチという意味では同じであったことを確認する。

\subsubsection{コマ主題細目}
\begin{description}
    \item[因子分析とテスト理論]単因子モデルの特殊事例として，学力テストの例を考える。学力テストの性質から，因子構造よりも因子得点に注目するという強調点の違いはあるが，因子分析モデルの一環として捉えることを強調する。
    \item[通過率と累積正規分布] 学力テストの分析例として，通過率の計算から累積正規分布へとつなげる。累積正規分布をそのまま確率モデルに繋げてもよいが，ロジスティック曲線を使う方が便利であることへつなげる。一般化線形モデルの文脈で考えれば，ロジスティック回帰分析をしていることにもなる。
    \item[項目母数の特徴]ロジスティック曲線を導入することで，関数の変形がたやすくなった。パラメータを1，2，3と増やすことでどこがどのように変化し，それぞれどういう意味があるのかについて確認する。
    \item[被験者母数の特徴]被験者母数の推定方法について確認する。手元の情報を全て使い切って推定する，完全情報最尤推定法であり，古典的テスト理論の欠点を克服できることを理解する。
\end{description}
\subsubsection{キーワード}
\begin{itemize}
\item 通過率
\item ロジスティックモデル
\item 被験者母数の推定について
\end{itemize}
\subsection{授業情報}
\paragraph{コマの展開方法}
講義
\subsubsection{予習・復習課題}
\paragraph{予習} テストの前提となる標準正規分布について復習しておく。特にRを使って出力できる確率密度，確率点，累積確率など手を動かして予習しておくと良い。
\paragraph{復習} 適当なグラフ描画ツール(Rでよい)をつかって，ロジスティックモデルを描写し，項目母数をどのように変えるとどのように曲線の形が変わるかを確認してみよう。
